\documentclass[10pt]{article}
\usepackage[margin=0.75in]{geometry}
\usepackage{amsmath, amssymb, amsthm}
\usepackage{microtype}
\usepackage{enumitem}
\usepackage{fancyhdr}
\usepackage{titlesec}
\usepackage{tcolorbox}
\usepackage{booktabs}

\linespread{1.08}
\setlength{\parskip}{0.4ex plus 0.2ex minus 0.1ex}

\titleformat{\section}{\large\bfseries}{\thesection.}{0.5em}{}
\titleformat{\subsection}{\normalsize\bfseries}{\thesubsection}{0.5em}{}
\titlespacing*{\section}{0pt}{1.5ex}{1ex}
\titlespacing*{\subsection}{0pt}{1ex}{0.5ex}

\setlist{itemsep=1pt, topsep=3pt, parsep=1pt, leftmargin=1.5em}

\pagestyle{fancy}
\fancyhf{}
\fancyhead[L]{\small EECS 127}
\fancyhead[R]{\small Discussion 7}
\fancyfoot[C]{\small\thepage}
\renewcommand{\headrulewidth}{0.4pt}

\newcommand{\RR}{\mathbb{R}}
\newcommand{\NN}{\mathcal{N}}
\newcommand{\Range}{\mathcal{R}}
\DeclareMathOperator{\rank}{rank}
\DeclareMathOperator{\trace}{tr}
\DeclareMathOperator{\dom}{dom}
\DeclareMathOperator{\epi}{epi}

\begin{document}

\begin{center}
    {\large\bfseries EECS 127/227AT \quad Optimization Models in Engineering \quad UC Berkeley \quad Spring 2026}\\[6pt]
    {\LARGE\bfseries Discussion 7 -- Solutions}
\end{center}

\section{Convexity of Functions}

\underline{Definition.} A function $f\colon \RR^n \to \RR$ is convex if $\dom(f)$ is a convex set and if for all $\vec{x}, \vec{y} \in \dom(f)$ and $\theta \in [0,1]$, we have,
\begin{equation}
    f(\theta \vec{x} + (1-\theta)\vec{y}) \leq \theta f(\vec{x}) + (1-\theta) f(\vec{y}).
\end{equation}
The function $f$ is strictly convex if the inequality is strict.

\underline{Definition.} A function $f\colon \RR^n \to \RR$ is concave if $\dom(f)$ is a convex set and if for all $\vec{x}, \vec{y} \in \dom(f)$ and $\theta$ with $0 \leq \theta \leq 1$, we have,
\begin{equation}
    f(\theta \vec{x} + (1-\theta)\vec{y}) \geq \theta f(\vec{x}) + (1-\theta) f(\vec{y}).
\end{equation}
The function $f$ is strictly concave if the inequality is strict.

\underline{Definition.} The epigraph $\epi(f)$ of a function $f\colon \Omega \to \RR$ is the set $\{(\vec{x}, t) \in \Omega \times \RR \mid t \geq f(x)\}$.

\textbf{Property.} A function $f$ is concave if and only if $-f$ is convex. An affine function is both convex and concave.

\textbf{Property: Jensen's inequality.} The inequality in Equation~(1) is known as \textbf{Jensen's Inequality}. This can be extended to convex combinations of more than one point. If $f$ is convex, and $\vec{x}_1, \vec{x}_2, \ldots, \vec{x}_k \in \dom(f)$, and $\theta_1, \theta_2, \ldots, \theta_k \geq 0$ with $\sum_{i=1}^{k} \theta_i = 1$ then,
\begin{equation}
    f(\theta_1 \vec{x}_1 + \theta_2 \vec{x}_2 + \cdots + \theta_k \vec{x}_k) \leq \theta_1 f(\vec{x}_1) + \theta_2 f(\vec{x}_2) + \cdots + \theta_k f(\vec{x}_k).
\end{equation}

\textbf{Property: first order condition.} Suppose $f$ is differentiable. Then $f$ is convex if and only if $\dom(f)$ is convex and
\begin{equation}
    f(\vec{y}) \geq f(\vec{x}) + \nabla f(\vec{x})^\top (\vec{y} - \vec{x}),
\end{equation}
for all $\vec{x}, \vec{y} \in \dom(f)$.

\textbf{Property: Second order condition.} Suppose $f$ is twice differentiable. Then $f$ is convex if and only if, $\dom(f)$ is convex and the Hessian of $f$, $\nabla^2 f(\vec{x})$, is positive semi-definite for all $\vec{x} \in \dom(f)$.

\textbf{Property: Convexity of the epigraph.} A function $f\colon \Omega \to \RR$, where $\Omega$ is a vector space, is convex if and only if its epigraph is convex, where linear combinations of elements of $\Omega \times \RR$ are defined in the obvious way.

\begin{enumerate}[label=(\alph*)]
    \item \textbf{Restriction to a line.}

    Show that a function $f$ is convex if and only if for all $\vec{x} \in \dom(f)$ and all $\vec{v}$, the function $g\colon \dom(g) \to \RR$ given by $g(t) = f(\vec{x} + t\vec{v})$ is convex for $\dom(g) = \{t \in \RR \mid \vec{x} + t\vec{v} \in \dom(f)\}$.

    \textbf{Solution:} \underline{In the first direction}: assume $f$ is convex and consider $\vec{x} \in \dom(f)$, $\vec{v}$ and the function $g\colon \dom(g) \to \RR$ given by $g(t) = f(\vec{x} + t\vec{v})$ where $\dom(g) = \{t \in \RR \mid \vec{x} + t\vec{v} \in \dom(f)\}$.

    Because $f$ is convex, $\dom(f)$ is convex, therefore $\dom(g)$ is also convex. For $t_1, t_2 \in \dom(g)$ and $\lambda \in [0,1]$:
    \begin{align}
        g(\lambda t_1 + (1-\lambda) t_2) &= f(\vec{x} + (\lambda t_1 + (1-\lambda) t_2)\vec{v}) \\
        &= f(\lambda(\vec{x} + t_1 \vec{v}) + (1-\lambda)(\vec{x} + t_2 \vec{v})) \\
        &\leq \lambda f(\vec{x} + t_1 \vec{v}) + (1-\lambda) f(\vec{x} + t_2 \vec{v}) \\
        &= \lambda g(t_1) + (1-\lambda) g(t_2)
    \end{align}
    Therefore $g$ is convex.

    \underline{In the other direction}: Consider $\vec{x}_1, \vec{x}_2 \in \dom(f)$ and $\lambda \in [0,1]$. Define $g\colon t \to f(\vec{x}_2 + t(\vec{x}_1 - \vec{x}_2))$. $g$ is convex and $0 \in \dom(g)$ and $1 \in \dom(g)$, so $[0,1] \in \dom(g)$. Therefore $\lambda \vec{x}_1 + (1-\lambda)\vec{x}_2 \in \dom(f)$ and $\dom(f)$ is convex.

    Because $g$ is convex:
    \begin{align}
        g(\lambda 1 + (1-\lambda)0) = g(\lambda) &\leq \lambda g(1) + (1-\lambda) g(0) \\
        f(\vec{x}_2 + \lambda(\vec{x}_1 - \vec{x}_2)) &\leq \lambda f(\vec{x}_2 + 1(\vec{x}_1 - \vec{x}_2)) + (1-\lambda) f(\vec{x}_2 + 0(\vec{x}_1 - \vec{x}_2)) \\
        f(\lambda \vec{x}_1 + (1-\lambda)\vec{x}_2) &\leq \lambda f(\vec{x}_1) + (1-\lambda) f(\vec{x}_2)
    \end{align}
    Therefore $f$ is convex.

    \item \textbf{Point-wise maximum.}

    Show that if $f_1$ and $f_2$ are convex functions then their pointwise maximum $f$, defined by
    \begin{equation}
        f(\vec{x}) = \max(f_1(\vec{x}), f_2(\vec{x})),
    \end{equation}
    with $\dom(f) = \dom(f_1) \cap \dom(f_2)$, is also convex.

    \textbf{Solution:}

    \underline{\textbf{Solution 1}}

    Because $f_1$ and $f_2$ are convex, then $\dom(f_1)$ and $\dom(f_2)$ are convex sets. Because convexity of sets is preserved under intersection, $\dom(f) = \dom(f_1) \cap \dom(f_2)$ is also convex.
    \begin{align}
        \epi(f) &= \{(\vec{x}, t) \mid \vec{x} \in \dom(f),\, f(\vec{x}) \leq t\} \\
        &= \{(\vec{x}, t) \mid \vec{x} \in \dom(f),\, \max(f_1(\vec{x}), f_2(\vec{x})) \leq t\} \\
        &= \{(\vec{x}, t) \mid \vec{x} \in \dom(f_1) \cap \dom(f_2),\, f_1(\vec{x}) \leq t \text{ and } f_2(\vec{x}) \leq t\} \\
        &= \{(\vec{x}, t) \mid \vec{x} \in \dom(f_1),\, f_1(\vec{x}) \leq t\} \cap \{(\vec{x}, t) \mid \vec{x} \in \dom(f_2),\, f_2(\vec{x}) \leq t\} \\
        &= \epi(f_1) \cap \epi(f_2)
    \end{align}
    Because $f_1$ and $f_2$ are convex, then $\epi(f_1)$ and $\epi(f_2)$ are convex. Because convexity of sets is preserved under intersection, $\epi(f)$ is convex. Because of the equivalence between the convexity of functions and the convexity of their epigraphs, $f$ is convex.

    \underline{\textbf{Solution 2}}

    In general, $\max(a+b, c+d) \leq \max(a,c) + \max(b,d)$. To show this formally, you can case-split on whether $a \geq c$ and whether $b \geq d$.

    This justifies the following steps:
    \begin{align}
        f(\lambda x_1 + (1-\lambda)x_2) &= \max(f_1(\lambda x_1 + (1-\lambda)x_2),\, f_2(\lambda x_1 + (1-\lambda)x_2)) \\
        &\leq \max(\lambda f_1(x_1) + (1-\lambda)f_1(x_2),\, \lambda f_2(x_1) + (1-\lambda)f_2(x_2)) \\
        &\leq \max(\lambda f_1(x_1),\, \lambda f_2(x_1)) + \max((1-\lambda)f_1(x_2),\, (1-\lambda)f_2(x_2)) \\
        &= \lambda \max(f_1(x_1),\, f_2(x_1)) + (1-\lambda)\max(f_1(x_2),\, f_2(x_2)) \\
        &= \lambda f(x_1) + (1-\lambda) f(x_2)
    \end{align}

    \underline{\textbf{Solution 3}}

    The idea is to use case-splitting so that we can simplify away a $\max$ operation in each case.

    For a given $x_1, x_2, \lambda$, either $f(\lambda x_1 + (1-\lambda)x_2) = f_1(\lambda x_1 + (1-\lambda)x_2)$ or $f(\lambda x_1 + (1-\lambda)x_2) = f_2(\lambda x_1 + (1-\lambda)x_2)$.

    In the first case,
    \begin{align}
        f(\lambda x_1 + (1-\lambda)x_2) &= f_1(\lambda x_1 + (1-\lambda)x_2) \\
        &\leq \lambda f_1(x_1) + (1-\lambda) f_1(x_2) \\
        &\leq \lambda \max(f_1(\vec{x}_1), f_2(\vec{x}_1)) + (1-\lambda)\max(f_1(\vec{x}_2), f_2(\vec{x}_2))
    \end{align}
    The other case is similar.
\end{enumerate}

\section{Convexity of Sets}

\underline{Definition.} A set $C$ is convex if and only if the line segment between any two points in $C$ lies in $C$:
\begin{equation}
    C \text{ is convex} \iff \forall\, \vec{x}_1, \vec{x}_2 \in C,\; \forall\, \theta \in [0,1],\; \theta \vec{x}_1 + (1-\theta)\vec{x}_2 \in C
\end{equation}

\begin{enumerate}[label=(\alph*)]
    \item Show that the \textbf{intersection of convex sets is convex}:
    \begin{equation}
        C_1, C_2 \text{ are convex} \implies C = C_1 \cap C_2 \text{ is convex}
    \end{equation}

    \textbf{Solution:} Consider $\vec{x}_1, \vec{x}_2 \in C$ and $\theta \in [0,1]$. Then $\vec{x}_1, \vec{x}_2 \in C_1$ and $\vec{x}_1, \vec{x}_2 \in C_2$. Since $C_1$ and $C_2$ are convex we have, $\theta \vec{x}_1 + (1-\theta)\vec{x}_2 \in C_1$ and $\theta \vec{x}_1 + (1-\theta)\vec{x}_2 \in C_2$, which implies $\theta \vec{x}_1 + (1-\theta)\vec{x}_2 \in C$.
\end{enumerate}

\end{document}
